\documentclass[12pt,a4paper]{article}

\usepackage[spanish,es-nodecimaldot]{babel}
\usepackage[T1]{fontenc}
\usepackage[utf8]{inputenc}
\usepackage{lmodern}
\usepackage{geometry}
\usepackage{setspace}
\usepackage{booktabs}
\usepackage{array}
\usepackage{longtable}
\usepackage{graphicx}
\usepackage{float}
\usepackage{hyperref}
\usepackage{caption}

\geometry{margin=2.5cm}
\onehalfspacing
\hypersetup{colorlinks=true, linkcolor=black, urlcolor=blue, citecolor=black}

\title{\textbf{Primer Proyecto de Inteligencia Artificial 2026}\\
Deteccion de Ciberacoso en Redes Sociales}
\author{Angel Esteban Esquit Hernández - 23221\\
Javier Eduardo España Pacheco - 23361\\
Roberto José Barreda Siekavizza - 23354}
\date{\today}

\begin{document}

\maketitle
\tableofcontents
\newpage

\section{Resumen Ejecutivo}
Este informe presenta el desarrollo de un sistema de clasificacion supervisada para detectar ciberacoso en tweets.
El problema se formula como clasificación multiclase sobre seis categorías: \textit{age, ethnicity, gender, religion,
other\_cyberbullying, not\_cyberbullying}. El flujo completo incluye preprocesamiento NLP, vectorización de texto,
entrenamiento y Comparación de Modelos, Optimización de hiperparámetros, análisis de overfitting y pruebas de usuario.

\section{Presentación y Análisis de Datos}
\subsection{Fuente y estructura del dataset}
Se utilizó el archivo \texttt{data/raw/cyberbullying\_tweets.csv}, con 47,692 muestras y dos columnas principales:
\texttt{tweet\_text} y \texttt{cyberbullying\_type}.

\subsection{Análisis exploratorio}
Se verificó la distribución de clases y se observó un conjunto aproximadamente balanceado entre las seis categorías.
La figura de distribución se incluye en el notebook y puede exportarse al reporte final.

\subsection{Datos faltantes y atípicos}
No se detectaron faltantes relevantes en las columnas de trabajo para la corrida final.
El tratamiento aplicado fue eliminar filas con nulos en texto o etiqueta antes del entrenamiento.

\subsection{Transformaciones aplicadas}
Se aplicaron las siguientes transformaciones al texto:
\begin{itemize}
  \item Eliminación de URLs, menciones, hashtags, caracteres especiales y emojis.
  \item Normalización a minúsculas.
  \item Tokenización.
  \item Eliminación de stopwords.
  \item Lematización.
\end{itemize}
Luego se realizaron representaciones BoW y TF-IDF, usando TF-IDF como vectorización principal para los experimentos reportados.

\section{Partición de Datos y Metodología}
Se utilizó división estratificada train-test de 80\% y 20\% con semilla fija 42 para reproducibilidad.
La variable objetivo se codificó a valores numéricos para entrenamiento.

\section{Modelos Probados y Configuración Técnica}
\subsection{Modelos implementados en el proyecto}
\begin{itemize}
  \item Naive Bayes (MultinomialNB)
  \item Regresión Logística
  \item SVM lineal
  \item Gradient Boosting
  \item MLP (red neuronal densa)
\end{itemize}

\subsection{Modelos evaluados en la corrida reportada}
\begin{itemize}
  \item Naive Bayes
  \item Logistic Regression (tuned)
  \item Gradient Boosting (tuned)
\end{itemize}

\subsection{Optimización de hiperparámetros}
Se aplico \texttt{GridSearchCV} a dos modelos, con metrica objetivo F1 ponderado.

\begin{table}[H]
\centering
\caption{Mejores hiperparametros encontrados}
\begin{tabular}{p{4.2cm} p{3.0cm} p{7.0cm}}
\toprule
Modelo & Mejor CV F1 (weighted) & Mejores parametros \\
\midrule
Logistic Regression & 0.8122 & \{C: 1.0, class\_weight: balanced\} \\
Gradient Boosting & 0.8085 & \{learning\_rate: 0.1, max\_depth: 5, n\_estimators: 150\} \\
\bottomrule
\end{tabular}
\end{table}

\section{Resultados y Comparación de Modelos}
\subsection{Métricas en conjunto de prueba}

\begin{table}[H]
\centering
\caption{Comparación de Modelos en test}
\begin{tabular}{lcccc}
\toprule
Modelo & Accuracy & Precision & Recall & F1 \\
\midrule
Logistic Regression (tuned) & 0.8113 & 0.8166 & 0.8113 & 0.8134 \\
Gradient Boosting (tuned)   & 0.8125 & 0.8291 & 0.8125 & 0.8101 \\
Naive Bayes                 & 0.7568 & 0.7493 & 0.7568 & 0.7478 \\
\bottomrule
\end{tabular}
\end{table}

\subsection{Métricas por clase del mejor modelo}

\begin{table}[H]
\centering
\caption{Reporte por clase: Logistic Regression (tuned)}
\begin{tabular}{lccc}
\toprule
Clase & Precision & Recall & F1 \\
\midrule
age                 & 0.9528 & 0.9725 & 0.9625 \\
ethnicity           & 0.9736 & 0.9711 & 0.9723 \\
gender              & 0.8901 & 0.8226 & 0.8550 \\
not\_cyberbullying  & 0.5470 & 0.5419 & 0.5444 \\
other\_cyberbullying& 0.5628 & 0.6211 & 0.5905 \\
religion            & 0.9670 & 0.9337 & 0.9501 \\
\bottomrule
\end{tabular}
\end{table}

\subsection{Análisis train vs test (overfitting)}
La comparación train-test muestra brechas de aproximadamente 0.054 a 0.058 en F1 para los modelos evaluados,
lo cual sugiere sobreajuste moderado pero controlado.

\begin{table}[H]
\centering
\caption{F1 train vs test por modelo}
\begin{tabular}{lccc}
\toprule
Modelo & F1 Train & F1 Test & Gap \\
\midrule
Logistic Regression (tuned) & 0.8676 & 0.8134 & 0.0542 \\
Gradient Boosting (tuned)   & 0.8651 & 0.8101 & 0.0550 \\
Naive Bayes                 & 0.8059 & 0.7478 & 0.0581 \\
\bottomrule
\end{tabular}
\end{table}

\subsection{Figuras de soporte}
\begin{figure}[H]
    \centering
    \includegraphics[width=0.88\textwidth]{model_comparison.png}
    \caption{Comparación de métricas por modelo en test.}
\end{figure}

\begin{figure}[H]
    \centering
    \includegraphics[width=0.88\textwidth]{train_vs_test_gap.png}
    \caption{Brecha train-test por métrica y modelo.}
\end{figure}

\begin{figure}[H]
    \centering
    \includegraphics[width=0.82\textwidth]{confusion_matrix_best.png}
    \caption{Matriz de confusión del mejor modelo reportado.}
\end{figure}

\section{Justificación del Modelo Seleccionado}

\subsection{Análisis de Costo Computacional}
Se realizó un análisis integral del costo computacional considerando tiempo de entrenamiento, complejidad algorítmica
e inferencia. La Tabla~\ref{tab:cost} resume la comparación de los tres modelos evaluados:

\begin{table}[H]
\centering
\caption{Análisis de Costo computacional por modelo}
\label{tab:cost}
\begin{tabular}{p{2.5cm}cccc}
\toprule
Modelo & Complejidad & Tiempo Est. Train. (s) & Inf./Muestra (ms) & F1 Test \\
\midrule
Naive Bayes                 & $O(n \cdot d)$          & $\sim 35$   & $< 1$ & 0.7478 \\
Logistic Regression (tuned) & $O(n \cdot d \cdot iter)$& $\sim 120$  & 1--2  & 0.8134 \\
Gradient Boosting (tuned)   & $O(n \cdot d \cdot trees)$& $\sim 180$ & 3--5  & 0.8101 \\
\bottomrule
\end{tabular}
\end{table}

La ejecución completa del pipeline (preprocesamiento, vectorización, tuning con GridSearchCV cv=3, y evaluación)
tomó aproximadamente \textbf{442.15 segundos} con seed=42 en configuración de 1 núcleo.
Con paralelización (n\_jobs=-1), se logra reducir tiempo de GridSearchCV significativamente.

\subsection{Justificación de Selección: F1 vs Costo}
Se selecciona \textbf{Logistic Regression (tuned)} como modelo final por ofrecer el mejor balance entre:
\begin{itemize}
  \item \textbf{Desempeño}: F1 ponderado de 0.8134, apenas 0.0033 por debajo de Gradient Boosting (0.8101 en test).
  \item \textbf{Costo computacional}: Tiempo de entrenamiento $\sim$3.4x menor que Gradient Boosting (120s vs 180s).
  \item \textbf{Interpretabilidad}: Modelo lineal facilita explicabilidad de decisiones, crucial en aplicaciones de moderación.
  \item \textbf{Inferencia}: Predicción 2--3x más rápida que árboles (1--2 ms vs 3--5 ms por muestra).
  \item \textbf{Despliegue}: Tamaño de modelo y dependencias mínimas, ideal para producción.
\end{itemize}

Aunque Gradient Boosting obtiene accuracy levemente mayor (0.8125 vs 0.8113), Logistic Regression ofrece mejor
compromiso entre precisión y costo computacional para una aplicación de detección de ciberacoso en tiempo real.

En el análisis por clase, las categorías \textit{not\_cyberbullying} y \textit{other\_cyberbullying} son las más difíciles.
Como mejora futura se recomienda ampliar el espacio de n-gramas y evaluar características de caracteres
(\textit{char n-grams}) para aumentar separabilidad entre clases semánticamente cercanas.

\section{Pruebas de Usuario}
Se implementó una función que recibe un tweet, aplica preprocesamiento, vectoriza y retorna predicción.\\
Los resultados de prueba se registran en \texttt{results/user\_test\_predictions.txt}.

\section{Cumplimiento de Restricciones de la Consigna}
\begin{itemize}
  \item No se usaron modelos preentrenados.
  \item Todos los modelos fueron entrenados desde cero sobre el dataset del proyecto.
  \item Se incluyeron tablas de metricas, Matriz de confusión, comparacion train-test y ejemplos de clasificación.
\end{itemize}

\section{Conclusiones}
El proyecto cumple los componentes técnicos principales de la Opción B: pipeline NLP, vectorización,
comparación de arquitecturas, tuning de hiperparámetros y evaluación integral.

Como cierre académico, se recomienda acompañar este informe con las figuras y tablas finales validadas en la
última corrida reproducible, y adjuntar el paquete de código y recursos en formato comprimido,
con el PDF del informe entregado por separado.

\end{document}
